\section{Closure and Spine Clashes}
\label{sec:spines}

This section turns unsafety of a fixed word $w$ into a syntactic test on
the constraints. By \cref{lem:fixed-witness}, $w$ is unsafe on a side
exactly when the signed system of that side has a solution that is fixed by
the assignment involution $S$ and satisfies two prefix conditions along $w$.
We encode these conditions, together with their sign duals, by an extension
$\Psi^4_w$ with four chains of fresh variables, the spines
(\cref{sec:spine}). Two facts are needed before the spines, and they fix
the order of this section. First, a covariant system is satisfiable exactly
when it has no label clash, a condition on its finite closure
(\cref{sec:closure,sec:least}). Second, a clash-free system that is
flip-closed and sign coherent has a solution fixed by $S$, obtained by
polarity selection on its least shape (\cref{sec:selector}). Together they
reduce unsafety of $w$ to the absence of a label clash in $\Psi^4_w$
(\cref{thm:four-spine}). We then express the closure of $\Psi^4_w$ through
the closure of the original system and two sets $U_j,L_j$ of old variables
at each cut (\cref{sec:bridge-formula}). Classifying the clashes in these
terms (\cref{sec:events}) yields tests at single cuts and tests on suffix
comparisons. \Cref{sec:holes} packages these tests as a finite reader with
the rejected-path property.

\subsection{Closure and Supported Path Bounds}\label{sec:closure}
For this subsection, let $\chi$ be any finite covariant constraint system over a
finite variable set $\Vars$. Its \emph{closure} $D_\chi$ is the least reflexive and
transitive relation satisfying, for every $i\in\Dir$,

\begin{equation}\label{eq:g21}
\frac{f(\vec a)\le u\in\chi\qquad D_\chi(u,v)\qquad
 v\le f(\vec b)\in\chi}{D_\chi(a_i,b_i)}.
\end{equation}

Repeatedly adding reflexive, transitive, and decomposition consequences
computes $D_\chi$: there are only $|\Vars|^2$ pairs. Induction on a derivation
proves soundness, $D_\chi(u,v)\Rightarrow \sigma(u)\le \sigma(v)$ for every solution.
For the decomposition step use the constructor case of \eqref{eq:g4}.

Define \emph{supported upper and lower path bounds} by

\begin{equation}\label{eq:g22}
\begin{aligned}
 \Up(\eps,u,v)&\iff D_\chi(u,v),\\
 \Up(i\eta,u,v)&\iff
 \exists z,\vec b\ [D_\chi(u,z)\land z\le f(\vec b)\in\chi
 \land\Up(\eta,b_i,v)],\\
 \Lo(\eps,u,v)&\iff D_\chi(u,v),\\
 \Lo(i\eta,u,v)&\iff
 \exists\vec a,z\ [f(\vec a)\le z\in\chi\land D_\chi(z,v)
 \land\Lo(\eta,u,a_i)].
 \end{aligned}
\end{equation}

In every solution, $\Up(\eta,u,v)$ says that the subtree of $u$ at
$\eta$, when present, is below $v$; $\Lo(\eta,u,v)$ says that $u$ is
below the subtree of $v$ there. The variable arguments follow the order of
the inequality: the path is on the first variable for $\Up$ and on the
last for $\Lo$.
To state soundness without assuming the path exists, let
$\emb_b(\eta,t)$ embed $t$ at $\eta$, putting the constant
tree $b$ in every off-path child. Then

\begin{lemma}[Path rules]\label{lem:path-rules}
In every solution of $\chi$,
\begin{equation}\label{eq:g23}
\begin{aligned}
 \Up(\eta,u,v)&\Longrightarrow
 \sigma(u)\le\emb_{\top}(\eta,\sigma(v)),\\
 \Lo(\eta,u,v)&\Longrightarrow
 \emb_{\bot}(\eta,\sigma(u))\le \sigma(v).
 \end{aligned}
\end{equation}

Composition and factorization, including empty factors, satisfy
\begin{equation}\label{eq:path-factors}
\begin{aligned}
\Up(\eta\tau,u,v)&\iff\exists a\;[\Up(\eta,u,a)\land\Up(\tau,a,v)],\\
\Lo(\eta\tau,u,v)&\iff\exists a\;[\Lo(\tau,u,a)\land\Lo(\eta,a,v)].
\end{aligned}
\end{equation}
At an empty factor the path judgment retains the closure relation
$D_\chi$, rather than requiring equality of its endpoints.
For later use, cancellation of a common path gives

\begin{equation}\label{eq:g24}
\begin{aligned}
 \Lo(\eta,a,z)\land
 \Up(\eta\tau,z,b)
 &\Longrightarrow\Up(\tau,a,b),\\
 \Lo(\eta\tau,a,z)\land
 \Up(\eta,z,b)
 &\Longrightarrow\Lo(\tau,a,b).
 \end{aligned}
\end{equation}

\end{lemma}
\begin{proof}
Induction on \eqref{eq:g22}, using \eqref{eq:g21} and the extreme
off-path fillers, proves \eqref{eq:g23}. Induction on the length of
$\eta$ proves both directions of \eqref{eq:path-factors}; closure
transitivity absorbs the intermediate pairs at empty factors.
For cancellation, when $\eta=\eps$, this is transitivity. For $\eta=i\eta'$, the first
lower and upper constructor bounds $f(\vec a')\le z'$ and
$z''\le f(\vec b')$, with their intervening closure pairs,
give $D_\chi(a'_i,b'_i)$ by \eqref{eq:g21}; apply the induction hypothesis to $\eta'$.
This proves both formulas and supplies the claimed factor and cancellation
rules in all arities. \end{proof}

A signed variable set has a sign map $\operatorname{sg}$ and a
fixed-point-free variable flip $\flip$ complementing signs. A system is
\emph{flip-closed} if reversing inequalities, flipping variables, and
exchanging constants preserves its literals. It is \emph{sign coherent}
if every constructor child $a_i$ of root $z$ has sign
$\operatorname{sg}(z)\oplus c(i)$, in either inequality orientation.

Return now to $\psi=\phi^{\pm}$. The literals of $\psi$ are closed under reversing
inequalities, flipping all signs, and exchanging constants. Induction on
\eqref{eq:g21} and \eqref{eq:g22} therefore proves

\begin{equation}\label{eq:g25}
\begin{aligned}
 D_\psi(u,v)&\iff D_\psi(\flip v,\flip u),\\
 \Lo(\eta,v,\flip u)
 &\iff\Up(\eta,u,\flip v).
 \end{aligned}
\end{equation}

The signed translation is sign coherent by its defining table.
\begin{lemma}[Sign purity]\label{lem:sign-purity}
In a sign-coherent system, closure pairs have equal signs. In either
path judgment along $\eta$, the endpoint signs differ by $p(\eta)$.
\end{lemma}
\begin{proof}
Reflexivity and transitivity preserve equality of signs. Decomposition
compares children at the same position, so both signs toggle by $c(i)$.
Induction on \eqref{eq:g22} proves the path assertion.
\end{proof}
In particular, a path from a positive query variable along $\eta$ has
its other endpoint of sign $p(\eta)$. Both properties refer to the signs
only through sign differences, so they also hold for the
\emph{complementary sign map} $1\oplus\operatorname{sg}$, under which the
negative copies are the positive variables.

\subsection{Label Clashes and the Least Shape}\label{sec:least}
A supported lower root label at $z$ is $f$ if some
$f(\vec a)\le u\in\chi$ has $D_\chi(u,z)$, and is $\top$ if some
$u=\top\in\chi$ has $D_\chi(u,z)$. Dually, an upper root label is
$f$ or $\bot$ if it is supplied by $v\le f(\vec b)$ or
$v=\bot$, with $D_\chi(z,v)$. A \emph{label clash} is a supported pair
$(\top,\bot)$, $(\top,f)$, or $(f,\bot)$. Equivalently, it is a
lower literal and an upper literal with a closure pair between their roots,
where the root label of the lower literal exceeds that of the upper one.


Define the three support predicates by
\begin{equation}\label{eq:shape-supports}
\begin{aligned}
\ell_f(z,\eta)&\iff\exists v,\vec a\;[\Lo(\eta,v,z)\land f(\vec a)\le v\in\chi],\\
u_f(z,\eta)&\iff\exists v,\vec b\;[\Up(\eta,z,v)\land v\le f(\vec b)\in\chi],\\
u_\bot(z,\eta)&\iff\exists v\;[\Up(\eta,z,v)\land v=\bot\in\chi].
\end{aligned}
\end{equation}
We retain a constructor when lower and upper constructor support force
one. Otherwise, we choose an extreme label compatible with the upper
bounds. A prescribed top needs no separate branch: in a clash-free
system it excludes the upper supports that could select another label.
Define
\begin{equation}\label{eq:g27}
 h(z,\eta)=
 \begin{cases}
 f,&\ell_f(z,\eta)\land u_f(z,\eta),\\
 \bot,&\neg\ell_f(z,\eta)\land(u_\bot(z,\eta)\lor u_f(z,\eta)),\\
 \top,&\text{otherwise}.
 \end{cases}
\end{equation}
The \emph{least-shape assignment} $\sigma_{\mathrm{sh}}$ retains this
label exactly at paths whose proper prefixes all have label $f$.
The root is retained, and a child is retained exactly when its parent is
labeled $f$, so this defines trees even when their domains are infinite.
Here ``least'' refers to the set of constructor positions, as proved
below, rather than to order-leastness: for binary covariant $f(a,a)\le a$
the rule gives $\top$, while the least solution is the infinite constructor tree.

The construction makes the least constructor shape explicit; its
common shape under the assignment involution enables the polarity
selection in \cref{lem:selector}.

\begin{theorem}[Least shape and satisfiability]\label{thm:sat}
If $\chi$ has no label clash, $\sigma_{\mathrm{sh}}\models\chi$. Hence
\begin{equation}\label{eq:g26}
 \Sat(\chi)\quad\Longleftrightarrow\quad
 \chi\text{ has no label clash}.
\end{equation}
\end{theorem}
\begin{proof}
First, the label rule is monotone in lower constructor support and
antitone in both upper supports. More precisely, if lower support at
$(v,\eta)$ implies that at $(v',\eta')$, and each upper support at
$(v',\eta')$ implies its counterpart at $(v,\eta)$, then
$h(v,\eta)\le h(v',\eta')$. For $h(v',\eta')=\top$, this is
automatic. For $h(v',\eta')=f$, upper constructor support at $(v,\eta)$ gives
$h(v,\eta)\le f$. For $h(v',\eta')=\bot$, absence of lower constructor support
transfers to $(v,\eta)$, as does the upper support causing bottom, so
$h(v,\eta)=\bot$.

For $f(\vec a)\le z$, the root supports a lower constructor, so
$f\le h(z,\eps)$. At a common child path $i\eta$, lower support at
$(a_i,\eta)$ propagates to $(z,i\eta)$ by \eqref{eq:g22}, and upper
support propagates in the reverse direction by \eqref{eq:g24}.
The monotonicity just proved gives the required comparison. For
$z\le f(\vec b)$, upper constructor support gives the root comparison;
cancellation transports lower support from $(z,i\eta)$ to $(b_i,\eta)$,
and the path rule transports upper support back. Thus this inequality
also holds at every common path. Pruning changes no label on common paths.
A prescribed bottom has upper bottom support and cannot have lower
constructor support without a clash, so the rule selects bottom.
A prescribed top excludes both upper supports and hence selects top.
This verifies every literal. Conversely, closure soundness excludes a
solution in the presence of an incompatible pair of supported labels.
\end{proof}

\begin{lemma}[Minimal constructor positions]\label{lem:minimal-shape}
If $\chi$ has no label clash, every constructor position of
$\sigma_{\mathrm{sh}}(z)$ is a constructor position of every solution's
value at $z$.
\end{lemma}
\begin{proof}
Induct along a retained path. At each constructor position $\eta$,
\eqref{eq:g27} supplies both constructor supports. By \eqref{eq:g23},
together with their supporting literals, they bound the solution
between $\emb_\bot(\eta,f(\sigma(a_i))_i)$ and
$\emb_\top(\eta,f(\sigma(b_i))_i)$ for suitable children. Both
embeddings have label $f$ at $\eta$. At the root the solution's label is therefore $f$; at subsequent
positions the induction hypothesis ensures the path exists, and the
same sandwich forces its label to be $f$. Since \cref{thm:sat} makes
$\sigma_{\mathrm{sh}}$ a solution, its constructor set is least by inclusion.
\end{proof}

\subsection{Polarity Selection}\label{sec:selector}
Use the assignment involution $S$ of \cref{lem:signed} on any signed
variable set. The signed translation is flip-closed and sign coherent, for
its own sign map and for the complementary one.
The proofs of \eqref{eq:g25} and sign purity apply to every such system;
flip closure also implies $\sigma\models\chi\Rightarrow S\sigma\models\chi$
by evaluating each flipped literal and reversing its tree inequality.

\begin{lemma}[Common shape and selection]\label{lem:selector}
For a flip-closed system, $\sigma_{\mathrm{sh}}$ and
$S\sigma_{\mathrm{sh}}$ have the same domains and constructor positions.
If it is also sign coherent and has no label clash, the assignment
\begin{equation}\label{eq:g60}
 \sigma_{\mathrm{sel}}(z)(\eta)=
 \begin{cases}
 \sigma_{\mathrm{sh}}(z)(\eta),&\operatorname{sg}(z)\oplus p(\eta)=0,\\
 (S\sigma_{\mathrm{sh}})(z)(\eta),&\operatorname{sg}(z)\oplus p(\eta)=1
 \end{cases}
\end{equation}
on this common domain satisfies $\chi$, is fixed by $S$, and preserves
that domain and its constructor positions exactly.
\end{lemma}
\begin{proof}
Variable flip exchanges lower and upper constructor support at $(z,\eta)$ and
$(\flip z,\eta)$. Thus \eqref{eq:g27} labels one by $f$ iff it labels the
other by $f$. Their proper-prefix pruning predicates coincide. Tree duality
preserves domains and constructor labels, giving the common shape.
Both assignments solve $\chi$ by \cref{thm:sat} and flip closure.
Selecting between them retains exactly their constructor positions,
so it gives a well-formed tree on their common domain.

At a constructor-literal root either input has the required comparison
with $f$. At a path $i\eta$, selection on the root side uses the bit
$\operatorname{sg}(z)\oplus c(i)\oplus p(\eta)$, the same bit used on
the child side at $\eta$, by coherence. Both labels therefore come from
one satisfying assignment. This proves either orientation of the literal.
Both inputs have the prescribed leaf at a constant literal.
Finally, at $\flip z$ the choice bit is complemented. If the bit at $z$
is zero, the selected label at $\flip z$ is the dual of the first input
at $z$; if it is one, the selected label there is the first input at
$\flip z$. Applying tree duality in either case gives the selected label
at $z$. Common domains prove the same equality at absent paths, so
$S\sigma_{\mathrm{sel}}=\sigma_{\mathrm{sel}}$.
\end{proof}

In particular, normalization in one direction and selection followed by
decoding in the other prove
\begin{equation}\label{eq:g19}
 \Sat_c(\phi)\quad\Longleftrightarrow\quad\Sat(\phi^\pm).
\end{equation}
Together with the finite closure test, this decides satisfiability.
The decoded solution is explicit. For $\psi=\phi^\pm$ without a label clash,
decoding $\sigma_{\mathrm{sel}}$ by \cref{lem:signed} gives the solution
\begin{equation}\label{eq:variance-shape}
 \sigma_{\mathrm{var}}(v)(\eta)=\sigma_{\mathrm{sh}}\bigl(v^{p(\eta)}\bigr)(\eta)
\end{equation}
of $\phi$, the \emph{variance shape}: at each path it reads the least shape
of the copy whose sign is the polarity of the path. Indeed
$\sigma_{\mathrm{var}}(v)=H_0(\sigma_{\mathrm{sel}}(v^+))$; when
$p(\eta)=0$ neither $H_0$ nor \eqref{eq:g60} changes the label of
$\sigma_{\mathrm{sh}}(v^+)$ at $\eta$, and when $p(\eta)=1$ the label is that
of $\sigma_{\mathrm{sh}}(v^-)$ dualized twice, by $S$ and by $H_0$.

\begin{example}[Selecting on the least shape]\label{ex:running-selection}
The running signed system of \cref{ex:running-signed} has identity closure
and no label clash. Both constructor bounds occur at $x^+$ and $x^-$;
their children have no constructor or constant bounds. Hence
\eqref{eq:g27} gives
\[
\begin{array}{c|cccc}
 v&x&y&a&b\\\hline
 \sigma_{\mathrm{sh}}(v^+)&f(\Top,\Top)&\Bot&\Top&\Top\\
 \sigma_{\mathrm{sh}}(v^-)&f(\Top,\Top)&\Top&\Top&\Top
\end{array}
\]
This solution is not fixed: at $x^+$ its image under $S$ is $f(\Bot,\Bot)$.
The domains nevertheless agree. Selection at $x^+$ takes the domain
leaf from $S\sigma_{\mathrm{sh}}$ and the codomain leaf from
$\sigma_{\mathrm{sh}}$, as in \cref{fig:polarity-selection}. In full, it gives
\[
\begin{array}{c|cccc}
 v&x&y&a&b\\\hline
 \sigma_{\mathrm{sel}}(v^+)&f(\Bot,\Top)&\Bot&\Top&\Top\\
 \sigma_{\mathrm{sel}}(v^-)&f(\Top,\Bot)&\Top&\Bot&\Bot
\end{array}
\]
Decoding gives another solution of $\phi$: $x=f(\Top,\Top)$,
$y=\Bot$, and $a=b=\Top$.
\end{example}

\begin{figure}[t]
\centering
\begin{tikzpicture}[x=10mm,y=8mm,
    every node/.style={font=\small,inner sep=2pt}]
  \foreach \name/\pos/\leftlab/\rightlab in
    {h/0/\Top/\Top,d/3.7/\Bot/\Bot,s/7.4/\Bot/\Top} {
    \node (\name) at (\pos,0) {$f$};
    \node (\name0) at (\pos-0.8,-1.1) {$\leftlab$};
    \node (\name1) at (\pos+0.8,-1.1) {$\rightlab$};
    \draw (\name) -- (\name0);
    \draw (\name) -- (\name1);
  }
  \node[above=2mm of h] {$\sigma_{\mathrm{sh}}(x^+)$};
  \node[above=2mm of d] {$(S\sigma_{\mathrm{sh}})(x^+)$};
  \node[above=2mm of s] {$\sigma_{\mathrm{sel}}(x^+)$};
  \node at (1.85,-0.6) {and};
  \node at (5.55,-0.6) {$\longmapsto$};
  \node[below=1mm of s0,font=\scriptsize,align=center]
    {bit $1$\\from $S\sigma_{\mathrm{sh}}$};
  \node[below=1mm of s1,font=\scriptsize,align=center]
    {bit $0$\\from $\sigma_{\mathrm{sh}}$};
\end{tikzpicture}
\caption{Selection in \cref{ex:running-selection}. All three trees have the
same shape. At the positive root the choice bit is $0$; at children $0,1$
it is respectively $1,0$, yielding $f(\Bot,\Top)$.}
\Description{Three trees with a constructor root and two leaves. The least
shape has two top leaves, its image under the assignment involution has two bottom leaves, and the selected
tree takes the left bottom leaf from the dual and the right top leaf from the
least shape. The choice bits are one at the left child and zero at the right.}
\label{fig:polarity-selection}
\end{figure}


\subsection{The Four-Spine Extension}\label{sec:spine}
Throughout \cref{sec:spine,sec:bridge-formula,sec:events}, $\psi$ may be any
flip-closed, sign-coherent system on a signed variable set, and $x^+,y^+$
any two positive variables, with $x^-=\flip x^+$ and $y^-=\flip y^+$.
In particular, this includes both sides of \cref{sec:fixed-witness}: the
left side is $\phi^{\pm}$ with the roots $(x^+,y^+)$, and the right side is
$\phi^{\pm}$ under the complementary sign map, for which its roots
$(X_{\mathsf r},Y_{\mathsf r})=(y^-,x^-)$ are positive. The assignment
involution $S$, and hence the notion of a fixed solution, depends only on
the variable flip and is the same for both sign maps.

Fix $w=w[0]\cdots w[m-1]$ and $p_j=p(w[:j])$.
To enforce the two fixed-witness conditions and their sign duals, we seek
\begin{equation}\label{eq:g29}
\begin{aligned}
 \emb_\bot(w,\top)&\le\sigma(x^+),&
 \sigma(x^-)&\le\emb_\top(w,\bot),\\
 \sigma(y^+)&\le\emb_\top(w,f(\top)_i),&
 \emb_\bot(w,f(\bot)_i)&\le\sigma(y^-).
\end{aligned}
\end{equation}
The four families flatten these bounds along $w$:
\begin{center}
\small
\begin{tabular}{@{}llll@{}}
\toprule
Spine & Bound direction & Intended condition along $w$ & Sign-dual partner\\
\midrule
$P$ & lower & $x^+$ has a top prefix & $Q$\\
$Q$ & upper & $x^-$ has a bottom prefix & $P$\\
$Y$ & upper & $y^+$ has no top prefix & $W$\\
$W$ & lower & $y^-$ has no bottom prefix & $Y$\\
\bottomrule
\end{tabular}
\end{center}
The families $P$ and $Q$ start at the two signed copies $x^+$ and $x^-$
of the query root $x$, and $Y$ and $W$ at $y^+$ and $y^-$. Thus $Q$ and
$W$ are the sign duals of $P$ and $Y$: flipping signs exchanges $P$ with
$Q$ and $Y$ with $W$.
The spines $P$ and $Y$ express the fixed-witness conditions. Their
sign-dual partners make the extension flip-closed, so selection yields a
fixed solution.

Set
$P_0=x^+$, $Q_0=x^-$, $Y_0=y^+$, $W_0=y^-$.
For $1\le j\le m$ introduce fresh base variables $x_j,y_j$, which carry the
spines of $x$ and $y$ at cut $j$, and name
\[
 P_j=x_j^{p_j},\quad Q_j=x_j^{1\oplus p_j},\qquad
 Y_j=y_j^{p_j},\quad W_j=y_j^{1\oplus p_j}.
\]
Two further base variables $o_0,o_1$, for the off-path fillers, supply four
signed fillers
$B^0=o_0^0$, $T^1=o_0^1$, $T^0=o_1^0$, $B^1=o_1^1$.
Thus $\flip B^q=T^{1-q}$, $\flip P_j=Q_j$, and $\flip Y_j=W_j$.
The variables with positive index and the fillers are fresh. The cut-zero
names $P_0=x^+$, $Q_0=x^-$, $Y_0=y^+$, $W_0=y^-$ are old variables;
$P_0=Y_0$ and $Q_0=W_0$ when $x=y$.
There are $2k+4m+4$ signed variables.
For $R\in\{P,W\}$ and $Z\in\{Q,Y\}$ put
\[
 a^{R,j}_i=\begin{cases}R_{j+1}&i=w[j],\\B^{\operatorname{sg}(R_j)\oplus c(i)}&i\ne w[j],\end{cases}
 \qquad
 b^{Z,j}_i=\begin{cases}Z_{j+1}&i=w[j],\\T^{\operatorname{sg}(Z_j)\oplus c(i)}&i\ne w[j].\end{cases}
\]
The extension $\Psi^4_w$ consists of $\psi$ and
\begin{equation}\label{eq:g28}
\begin{array}{ll}
 f(a^{P,j}_i)_i\le P_j,&f(a^{W,j}_i)_i\le W_j\quad(0\le j<m),\\
 Q_j\le f(b^{Q,j}_i)_i,&Y_j\le f(b^{Y,j}_i)_i\quad(0\le j<m),\\
 P_m=\top,&Q_m=\bot,\\
 Y_m\le f(T^{p_m\oplus c(i)})_i,&f(B^{1\oplus p_m\oplus c(i)})_i\le W_m,\\
 B^q=\bot,&T^q=\top\quad(q=0,1).
\end{array}
\end{equation}
It is sign coherent: the distinguished child uses
$p_{j+1}=p_j\oplus c(w[j])$, and all other children, including terminal
children, have the prescribed sign by definition. It is flip-closed:
flip pairs the $P$ and $Q$ links, the $W$ and $Y$ links, the terminal
constants, the terminal constructor bounds, and the opposite-sign filler
constants. These identities also hold at cut zero and when $m=0$.

\Cref{fig:four-spines} displays the four parallel families of bounds;
the arrows indicate the distinguished child, not closure pairs.
\begin{figure}[t]
\centering
\begin{tikzpicture}[x=9mm,y=8mm,>=Stealth,
    every node/.style={font=\small,inner sep=2pt},
    filler/.style={font=\scriptsize,anchor=west}]
  \node[anchor=east,font=\scriptsize] at (-0.5,0) {$f(\ldots)\le P_j$};
  \node[anchor=east,font=\scriptsize] at (-0.5,-1.6) {$Q_j\le f(\ldots)$};
  \node[anchor=east,font=\scriptsize] at (-0.5,-3.2) {$Y_j\le f(\ldots)$};
  \node[anchor=east,font=\scriptsize] at (-0.5,-4.8) {$f(\ldots)\le W_j$};
  \node at (0,0.7) {cut $0$};
  \node at (1.8,0.7) {$j$};
  \node at (3.8,0.7) {$j+1$};
  \node at (6.1,0.7) {$m$};
  \foreach \spinerow/\spiney/\spineroot/\spinesign/\spinenextsign/\spinefiller in
    {P/0/x^+/p_j/p_{j+1}/B,
     Q/-1.6/x^-/1\oplus p_j/1\oplus p_{j+1}/T,
     Y/-3.2/y^+/p_j/p_{j+1}/T,
     W/-4.8/y^-/1\oplus p_j/1\oplus p_{j+1}/B} {
    \node (\spinerow0) at (0,\spiney) {$\spineroot$};
    \node (\spinerow j) at (1.8,\spiney) {$\underset{[\spinesign]}{\spinerow_j}$};
    \node (\spinerow next) at (3.8,\spiney) {$\underset{[\spinenextsign]}{\spinerow_{j+1}}$};
    \draw[densely dotted] (\spinerow0) -- (\spinerow j);
    \draw[->] (\spinerow j) -- node[above,font=\scriptsize] {$w[j]$} (\spinerow next);
    \draw (\spinerow j.south) -- ++(0,-0.45) -- ++(0.4,0)
      node[filler] {$\spinefiller^{\spinesign\oplus c(i)}\quad(i\ne w[j])$};
  }
  \node[anchor=west] (Pm) at (5.5,0) {$\underset{[p_m]}{P_m}=\Top$};
  \node[anchor=west] (Qm) at (5.5,-1.6) {$\underset{[1\oplus p_m]}{Q_m}=\Bot$};
  \node[anchor=west] (Ym) at (5.5,-3.2)
    {$\underset{[p_m]}{Y_m}\le f(T^{p_m\oplus c(i)})_i$};
  \node[anchor=west] (Wm) at (5.5,-4.8)
    {$f(B^{1\oplus p_m\oplus c(i)})_i\le\underset{[1\oplus p_m]}{W_m}$};
  \foreach \spinerow in {P,Q,Y,W} {\draw[densely dotted] (\spinerow next) -- (\spinerow m);}
\end{tikzpicture}
\caption{The four-spine extension along $w$, with $Q_j=\flip P_j$ and
$W_j=\flip Y_j$. Brackets give node signs; the root signs are $0,1,0,1$.
Each arrow reads the child $w[j]$. The $P,W$ links are lower constructor
bounds, and the $Q,Y$ links are upper bounds, with the displayed off-path
fillers. Dots abbreviate the other cuts. \Cref{ex:running-spines}
instantiates the construction at $w=0$.}
\Description{Four horizontal rows P, Q, Y, W run from their signed query
roots through cuts j and j plus one to cut m. Every step reads the letter at position j and
toggles the sign by its variance. P and W have bottom off-path fillers;
Q and Y have top fillers. The final P is top and Q is bottom. Y ends with
an upper constructor bound and W with a lower constructor bound.}
\label{fig:four-spines}
\end{figure}


\begin{example}[One contravariant step]\label{ex:running-spines}
For the running query and $w=0$, the root signs of $(P,Q,Y,W)$ are
$(0,1,0,1)$, and their signs at cut $1$ are $(1,0,1,0)$.
The new links and terminal bounds are exactly
\[
\begin{array}{ll}
 f(P_1,B^0)\le x^+,&x^-\le f(Q_1,T^1),\\
 y^+\le f(Y_1,T^0),&f(W_1,B^1)\le y^-,\\
 P_1=\Top,\quad Q_1=\Bot,&
 Y_1\le f(T^0,T^1),\quad f(B^1,B^0)\le W_1,
\end{array}
\]
together with $B^q=\Bot$ and $T^q=\Top$ for $q=0,1$.
Extend the fixed witness of \cref{ex:running-signed} by
\[
 \sigma(P_1)=\Top,\quad \sigma(Q_1)=\Bot,\qquad
 \sigma(Y_1)=f(\Top,\Top),\quad \sigma(W_1)=f(\Bot,\Bot).
\]
With the prescribed fillers this is a fixed solution of the extension,
which therefore has no label clash.
\end{example}

\begin{lemma}[Spine extension]\label{lem:spine}
An assignment $\sigma\models\psi$ extends to a solution of $\Psi^4_w$ iff
\begin{equation}\label{eq:g18}
\begin{array}{ll}
 \widehat{\sigma(x^+)}(w)=\top,&\widehat{\sigma(x^-)}(w)=\bot,\\
 \widehat{\sigma(y^+)}(w)\ne\top,&\widehat{\sigma(y^-)}(w)\ne\bot.
\end{array}
\end{equation}
\end{lemma}
\begin{proof}
Composing the links gives \eqref{eq:g29}.
The first bound forces a top at or before the final cut: an earlier bottom
violates the on-spine constructor bound, while reaching the end forces top.
Conversely, a top prefix satisfies the bound; all off-path bottom comparisons
are automatic. Duality proves the second equivalence. The third excludes
a top at every cut through the end, since the upper bound has constructors
there. A bottom ends comparison harmlessly; otherwise the terminal label
is at most $f$. Duality gives the fourth equivalence.
For sufficiency, assign, for each $1\le j\le m$,
\[
\begin{aligned}
\sigma(P_j)&=\emb_\bot(w[j:],\top),&
\sigma(Q_j)&=\emb_\top(w[j:],\bot),\\
\sigma(Y_j)&=\emb_\top(w[j:],f(\top)_i),&
\sigma(W_j)&=\emb_\bot(w[j:],f(\bot)_i).
\end{aligned}
\]
These are assignments to distinct fresh signed variables, so they impose
no extra condition on the old assignment, even when it is not fixed by
$S$. Assign each filler its prescribed constant.
For $m\ge1$, internal links and terminal literals hold with equality, and the root links
are the four bounds \eqref{eq:g29}; all newly assigned trees are finite,
so the extension stays in each of the three domains.
For $m=0$ the terminal literals themselves are those bounds at the old roots.
\end{proof}

\begin{theorem}[Four-spine unsafe-word criterion]\label{thm:four-spine}
On either side $\theta$, $w$ is $\theta$-unsafe iff the four-spine extension
of $\psi$ for the roots $(X_\theta,Y_\theta)$ has no label clash. No
satisfiability assumption is needed.
\end{theorem}
\begin{proof}
For the left side, an unsafe word gives a fixed witness by
\cref{lem:fixed-witness}. Fixedness supplies the opposite-sign conditions
in \eqref{eq:g18}, so \cref{lem:spine} extends it to a solution and excludes
a clash. Conversely, apply \cref{lem:selector} to the whole clash-free
extension \eqref{eq:g28}. The selected least shape is a fixed solution;
restriction to the old variables preserves satisfaction and fixedness,
and \cref{lem:spine} supplies all four prefix conditions. Its first and
third conditions give the fixed witness and hence unsafety.
For the right side, apply the same argument to $\psi$ with the roots
$(y^-,x^-)$, which are positive for the complementary sign map, and conclude
by \cref{lem:fixed-witness} for $\theta=\mathsf r$.
\end{proof}

\begin{lemma}[Separation]\label{lem:separation}
Partition the variables into old variables and the two fresh classes
\[
 \text{sources }\{P_j,W_j:1\le j\le m\}\cup\{B^0,B^1\},\qquad
 \text{sinks }\{Q_j,Y_j:1\le j\le m\}\cup\{T^0,T^1\}.
\]
Then

\begin{equation}\label{eq:g30}
\begin{array}{ll}
 D_{\Psi^4_w}(a,s),\ s\text{ a source}&\Longrightarrow a=s,\\
 D_{\Psi^4_w}(t,b),\ t\text{ a sink}&\Longrightarrow b=t,\\
 a,b\text{ old}:\quad D_{\Psi^4_w}(a,b)&\Longleftrightarrow D_\psi(a,b).
 \end{array}
\end{equation}

\end{lemma}

\begin{proof} Induct simultaneously on closure derivations.
An upper constructor literal has only old or sink children, never source
children; hence decomposition cannot have a source as its right endpoint.
Dually, a lower constructor literal has only old or source children, so
decomposition cannot have a sink as its left endpoint. This proves the
first two assertions, with transitivity collapsing the intervening identity.
For an old-to-old transitivity step, a fresh middle variable would be a source
or a sink and contradict these assertions. For an old-to-old decomposition
step, both bounding literals must be old: every new lower literal has only
source children, and every new upper literal has only sink children.
Apply induction to their old root pair. The other implication follows by
including the original derivation. Constants add no closure rule. \end{proof}

The classification here is by provenance of a literal, not the class of its
root: a new literal can be rooted at an old variable at cut zero.

\subsection{The Bridge Formula}\label{sec:bridge-formula}
Write $D=D_{\Psi^4_w}$ for the closure of the four-spine extension.
All path judgments in this subsection are in the base system $\psi$. Define

\begin{equation}\label{eq:g31}
U^R_j=\{v:\Up(w[:j],R_0,v)\},\qquad
 L^Z_j=\{v:\Lo(w[:j],v,Z_0)\},
\end{equation}

where $R\in\{P,W\}$ and $Z\in\{Q,Y\}$. Put $U_j=U^P_j$ and $L_j=L^Y_j$. By \eqref{eq:g25}, we have
$L^Q_j=\flip U_j$ and $U^W_j=\flip L_j$, where $\flip\mathcal U=\{\flip v:v\in\mathcal U\}$.
The sets $U^R_j$ are upward closed under $D_\psi$, and $L^Z_j$ are
downward closed. At cut zero,
$U^R_0=D_\psi\mathord\uparrow R_0$ and
$L^Z_0=D_\psi\mathord\downarrow Z_0$, where
$D_\psi\mathord\uparrow z=\{v:D_\psi(z,v)\}$ and
$D_\psi\mathord\downarrow z=\{v:D_\psi(v,z)\}$.

These sets already show why a clash can relate two different cuts. Take
cuts $s<e$ and $v\in U_e\cap L_s$. Where the subtrees exist, the path
bounds read
\[
 \sigma(x^+)|_{w[:e]}\le\sigma(v)\le\sigma(y^+)|_{w[:s]}.
\]
The upper embedding bound excludes a top leaf of $x^+$ before cut $e$.
If $x^+$ has a top leaf along $w$, its remaining path is therefore a
prefix of $w[e:]$, and the same bound forces a top prefix of $v$ there.
The lower embedding bound then forces a top leaf of
$y^+$ along $w[:s]w[e:]$. When $w[e:]\pre w[s:]$, this word is a prefix
of $w$, contradicting the target condition. The embedding formulation
handles an earlier top leaf of $y^+$ as well. In the extension, the old
variable $v$ links the $P$-spine at cut $e$ to the $Y$-spine at cut $s$,
and the equal suffixes carry this link down to the final cut. The bridge
formula describes all such links.

\begin{lemma}[Exact bridge formula]\label{lem:bridge}
For every old $v$, both lower names $R$, both upper names $Z$, and all $0\le j,j'\le m$,

\begin{equation}\label{eq:g32}
D_{\Psi^4_w}(R_j,v)\iff v\in U^R_j,
 \qquad D_{\Psi^4_w}(v,Z_{j'})\iff v\in L^Z_{j'},
\end{equation}

\begin{equation}\label{eq:g33}
\begin{split}
 D_{\Psi^4_w}(R_j,Z_{j'})\iff
 \exists \alpha\le j,\ \gamma\le j'\quad&
 U^R_\alpha\cap L^Z_\gamma\ne\varnothing,\\
 &j-\alpha=j'-\gamma,\qquad w[\alpha:j]=w[\gamma:j'].
 \end{split}
\end{equation}

\end{lemma}

The bridge consists of an old-variable intersection followed by parallel
descent through equal word factors (\cref{fig:bridge-descent}).
\begin{figure}[tbp]
\centering
\begin{tikzpicture}[x=12mm,y=9mm,>=Stealth,every node/.style={font=\small}]
\node (ra) at (0,0) {$R_\alpha$};
\node (v) at (3,0) {$v$ old};
\node (zg) at (6,0) {$Z_\gamma$};
\node (rj) at (0,-1.6) {$R_j$};
\node (zj) at (6,-1.6) {$Z_{j'}$};
\draw[->] (ra) -- node[above] {$D$} (v);
\draw[->] (v) -- node[above] {$D$} (zg);
\draw[->] (ra) -- node[left] {$w[\alpha:j]$} (rj);
\draw[->] (zg) -- node[right] {$w[\gamma:j']$} (zj);
\draw[->,thick] (rj) -- node[above] {$D$} (zj);
\node[font=\footnotesize] at (3,-2.3) {$w[\alpha:j]=w[\gamma:j']$};
\end{tikzpicture}
\caption{A bridge in \cref{lem:bridge}. The old variable belongs to
$U^R_\alpha\cap L^Z_\gamma$. Horizontal arrows are closure pairs;
vertical arrows descend through distinguished children along the same
factor. Decomposition propagates the top-row pair $D(R_\alpha,Z_\gamma)$ to $D(R_j,Z_{j'})$.
The common factor may be empty.}
\Description{A lower spine root connects through an old variable to an
upper spine root. Parallel descent along equal factors yields a closure
pair between the two later spine nodes.}
\label{fig:bridge-descent}
\end{figure}


\begin{proof}\emph{Right to left.} At zero, \eqref{eq:g32} follows from \cref{lem:separation}. Factor
an upper path along $w[:j]$ at its last letter. By induction the upper path along $w[:j-1]$
gives a closure pair from $R_{j-1}$ to the root of an old upper literal;
decompose against the $R$-link at position $w[j-1]$, then use the
terminal closure of the path judgment. This gives the right-to-left direction
of the first equivalence in \eqref{eq:g32}. The lower-path argument with a $Z$-link proves the second.
For \eqref{eq:g33}, choose $v\in U^R_\alpha\cap L^Z_\gamma$. By \eqref{eq:g32}, we have
$D(R_\alpha,v)$ and $D(v,Z_\gamma)$. Their composite is $D(R_\alpha,Z_\gamma)$.
Decompose the two spine links successively along the common factor
$w[\alpha:j]=w[\gamma:j']$. This gives $D(R_j,Z_{j'})$.

\emph{Left to right.} Strengthen the induction on a closure
derivation to four simultaneous obligations: old-to-old reflection in \eqref{eq:g30},
the two implications \eqref{eq:g32} for each lower or upper name at positive indices,
and \eqref{eq:g33} for each of the four name pairs when both indices are
positive. Pairs with a bottom filler on the left or a top filler on the
right carry no additional obligation; such pairs can occur.

For transitivity, an intermediate source forces its incoming pair to
be an identity, and an intermediate sink forces its outgoing pair to be
an identity, by \eqref{eq:g30}. An old intermediate variable joins the
two path memberships, producing the intersection in \eqref{eq:g33}.
Decomposition advances both spines along the same letter; empty factors
retain closure, so the invariant also covers cut-zero aliases.
The remaining rule cases and zero endpoints are detailed in
\cref{app:spine-proofs}.
\end{proof}

\subsection{Exact Clash Analysis}\label{sec:events}
Let $\mathsf{Bot}=\{v:v=\bot\in\psi\}$ and
$\mathsf{Top}=\{v:v=\top\in\psi\}$. Define

\begin{equation}\label{eq:g34}
\begin{aligned}
 \operatorname{bot}(U)&\iff U\cap\mathsf{Bot}\ne\varnothing,\\
 \operatorname{top}(L)&\iff L\cap\mathsf{Top}\ne\varnothing,\\
 \operatorname{refl}(U,L)&\iff U\cap L\ne\varnothing,\\
 \ready(U,L)&\iff
 \operatorname{bot}(U)\lor\operatorname{top}(L)\lor
 \operatorname{refl}(U,L),\\
 \childtest(U)&\iff
 \exists v\in U,\vec b\quad v\le f(\vec b)\in\psi.
 \end{aligned}
\end{equation}

At cut $j$, bot-readiness bounds $x^+$ by bottom along $w[:j]$,
excluding a top prefix along $w$. Top-readiness forces a top prefix of
$y^+$ by cut $j$. Reflexivity readiness bounds $x^+$ by $y^+$ at the
same cut, so any remaining top prefix transfers to $y^+$.
These readings use the embedding bounds \eqref{eq:g23}, and thus also
cover trees that terminate before the cut. The child test at the terminal
cut bounds $x^+$ by a constructor along all of $w$, excluding a top
prefix there. It does not depend on a next letter: it asks only whether
$U$ contains the root of an upper constructor bound.

The two-cut events concern the following pairs of cuts.

\begin{definition}[Comparison]\label{def:comparison}
A \emph{comparison} of a word $w$ of length $m$ is a pair of cuts $(s,e)$
satisfying
\begin{equation}\label{eq:g35}
0\le s<e\le m,\qquad w[e:]\preceq w[s:].
\end{equation}
Its \emph{period} is $\varpi=e-s$, and its \emph{root} is the factor
$w[s:e]$. A comparison is \emph{terminal} if $e=m$.
\end{definition}

\Cref{sec:holes} interprets comparisons as periodic suffixes and gives
examples. A comparison of $w$ is \emph{cross-admitted} if
$L_s\cap U_e\ne\varnothing$, and \emph{self-admitted}
if $U_s\cap\flip U_e\ne\varnothing$.

\begin{theorem}[Exact clash analysis]\label{thm:clashes}
For either side, the four-spine extension $\Psi^4_w$ has a label
clash iff $\psi$ has a label clash or at least one of the
following events occurs:

\begin{equation}\label{eq:g36}
\begin{array}{ll}
 \text{(a)}&\exists j\le m\quad\ready(U_j,L_j),\\
 \text{(b)}&\childtest(U_m),\\
 \text{(c)}&\exists s<e\le m\quad w[e:]\preceq w[s:]\ \land\
 [L_s\cap U_e\ne\varnothing\ \lor\ U_s\cap\flip U_e\ne\varnothing].
 \end{array}
\end{equation}

\end{theorem}

Readiness (a) rules out a fixed witness at a single cut, and the child event
(b) rules it out at the final cut. A cross event in (c) transfers a top
prefix from $x^+$ to $y^+$ across two cuts. A self event makes
the top-prefix requirement on $x^+$ incompatible with the bottom-prefix
requirement on $x^-$.

\begin{center}
\small
\begin{tabular}{@{}lcccc@{}}
\toprule
Lower bound $\backslash$ upper bound & old & $Q_{j'}$ & $Y_{j'}$ & $B^q$\\
\midrule
old & inherited clash & bot / child & top & impossible\\
$P_j$ & bot / child & self & refl / cross & impossible\\
$W_j$ & top & refl / cross & none ($f/f$) & impossible\\
$T^q$ & impossible & impossible & impossible & impossible\\
\bottomrule
\end{tabular}
\end{center}
The entries classify incompatible labels only; a closure pair between
compatible constructor labels is not a clash. Here bot, top, and refl
are the three readiness kinds, and child is the terminal event.

\begin{proof}Suppose first that a clash occurs. Express the clash by its two
bounding literals and a closure pair between their roots. The only relevant
new lower bounds are the $P_j$-bounds, with label $f$ for $j<m$
and $\top$ for $j=m$, the $W_j$-bounds, all labeled $f$ (including $j=m$), and the bounds $T^q=\top$. The relevant new
upper bounds are the $Q_{j'}$-bounds, labeled $f$ for $j'<m$ and
$\bot$ for $j'=m$, the $Y_{j'}$-bounds, all labeled $f$, and
$B^q=\bot$. The fillers cannot participate: a pair out of $T^q$, or into
$B^q$, is an identity by \eqref{eq:g30}, and neither has an opposite incompatible
bound. Classify the remaining bounds by whether their literals belong to
$\psi$ or were added in \eqref{eq:g28}.

If both are old, separation gives an inherited clash in $\psi$. If the lower bound is a $P_j$-bound and the upper bound
is old, an upper bottom yields $\operatorname{bot}(U_j)$ by \eqref{eq:g32}.
An upper constructor clashes only against $P_m=\top$, and yields
$\childtest(U_m)$.

The remaining old/new cases are their sign duals or give top-readiness; see \cref{app:spine-proofs}.

For a new $P$-bound against a new $Y$-bound, the $Y$-bound has label $f$,
so the lower endpoint must be $P_m$. Apply \eqref{eq:g33} to $D(P_m,Y_{j'})$. Its cuts satisfy
$m-\alpha=j'-\gamma$, so $\alpha\ge \gamma$. If $\alpha=\gamma$, the intersection
$U_\alpha\cap L_\alpha$ gives reflexivity readiness. If $\alpha>\gamma$, set
$s=\gamma,e=\alpha$. The factor equality reads
$w[e:m]=w[s:s+m-e]$; this is exactly \eqref{eq:g35}, and the intersection is
the cross admission $L_s\cap U_e\ne\varnothing$.

The $W$/$Q$ case is the sign dual of this cross case; $W$/$Y$ labels are compatible.

For a $P$-bound against a $Q$-bound, at least one endpoint is terminal:
$j=m$ or $j'=m$. The cuts $\alpha,\gamma$ given by \eqref{eq:g33} have
$U_\alpha\cap\flip U_\gamma\ne\varnothing$. Purity excludes $\alpha=\gamma$, since
$U_\alpha$ and $\flip U_\alpha$ have opposite signs. If $\alpha<\gamma$, the equation
$j-\alpha=j'-\gamma$ forces $j'>j$, hence $j'=m$; set $s=\alpha,e=\gamma$.
If $\alpha>\gamma$, it forces $j>j'$, hence $j=m$; set $s=\gamma,e=\alpha$ and
replace a common element $v$ by $\flip v$. In both cases the
equal factors give \eqref{eq:g35} and the intersection is
$U_s\cap\flip U_e\ne\varnothing$. These two orientations are both
needed. All other new/new pairs either have equal constructor labels or
cannot supply a lower and an upper incompatible label. This exhausts the
clashes, including literals rooted at old variables when $m=0$.

Conversely, an inherited clash persists. Readiness and child events
produce the clashes in the table, as detailed in \cref{app:spine-proofs}.
For a cross admission choose $v\in U_e\cap L_s$. The bridge yields
$D(P_e,Y_s)$. By \eqref{eq:g35}, descend $m-e$ matching letters to obtain
$D(P_m,Y_{s+m-e})$, a top/constructor clash with $P_m=\top$. For a self admission
choose $v\in U_s\cap\flip U_e=U_s\cap L^Q_e$. Descend from
$D(P_s,Q_e)$ to $D(P_{s+m-e},Q_m)$. Since $s<e$, the left index
is strictly less than $m$, so this is a constructor/bottom clash.
When $e=m$, no descent is needed, and the same clashes remain valid. \end{proof}

\begin{remark}[Special cases of the clash analysis]\label{rem:clash-cases}
\leavevmode
\begin{itemize}[leftmargin=*]
\item For satisfiable $\psi$, inherited clashes are absent, so the events
  (a)--(c) are exactly the extension clashes.
\item Purity shows that a self-admitted comparison must have
  $p(w[s:e])=1$: its two upper sets must have opposite signs. Thus
  self-admissions disappear in the fully covariant case, but cannot be
  discarded for general variance.
\item When $x=y$, reflexivity puts $x^+$ in $U_0\cap L_0$, so readiness at
  zero excludes every unsafe word.
\item When $m=0$, only readiness and the terminal child test remain; the
  zero-cut bridge formulas cover all aliases.
\item At arity zero, the constructor $f()$ still supplies a genuine
  terminal bound and is detected by the child test.
\item Child is tested only at the terminal cut; readiness is tested at
  every cut, including zero and the terminal cut.
\end{itemize}
\end{remark}

\begin{example}[A cross event]\label{ex:cross-event}
Replace the running system by $\phi=\{x\le f(a,y),\ y\le f(a,b)\}$,
keeping $c=(1,0)$ and the query $x\le y$. Its signed literals with positive
roots are $x^+\le f(a^-,y^+)$ and $y^+\le f(a^-,b^+)$; the other two are
their sign duals, and the closure is the identity. Along $w=11$, the upper sets are $\{x^+\},\{y^+\},\{b^+\}$ and
the lower sets are $\{y^+\},\varnothing,\varnothing$. No cut is ready, the
terminal child test fails, and there is no self event. The only event is
the comparison $(0,1)$, with $w[1:]=1\pre11=w[0:]$, which is cross-admitted
because $y^+\in U_1\cap L_0$. \Cref{fig:cross-event} shows the resulting
clash. The shared variable $y^+$ supplies the bridge $D(P_1,Y_0)$ between
two different cuts, and the matching suffix transports it one letter down
to $D(P_2,Y_1)$, where $P_2=\top$ meets the upper constructor bound of $Y_1$.
\end{example}
\begin{figure}[t]
\centering
\begin{tikzpicture}[x=12mm,y=8mm,>=Stealth,
    every node/.style={font=\small,inner sep=2pt}]
  \node (s) at (0,0) {$s=0$};
  \node (e) at (2,0) {$e=1$};
  \node (m) at (4,0) {$m=2$};
  \draw[->] (s) -- node[above] {$1$} (e);
  \draw[->] (e) -- node[above] {$1$} (m);
  \node[anchor=west] at (4.6,0) {$w[e:]=1\pre 11=w[s:]$};
  \node (p1) at (2,-1.4) {$P_1$};
  \node (bridge) at (1,-1.4) {$y^+$};
  \node (y0) at (0,-1.4) {$Y_0$};
  \draw[->] (p1) -- node[above,font=\scriptsize] {$D$} (bridge);
  \draw[->] (bridge) -- node[above,font=\scriptsize] {$D$} (y0);
  \node[anchor=west] at (4.6,-1.4) {$y^+\in U_1\cap L_0$};
  \node (p2) at (4,-3) {$P_2=\Top$};
  \node (y1) at (2,-3) {$Y_1\le f(\ldots)$};
  \draw[->] (p1) -- node[right,font=\scriptsize] {child $1$} (p2);
  \draw[->] (y0) -- node[left,font=\scriptsize] {child $1$} (y1);
  \draw[->,thick] (p2) -- node[above] {$D$}
    node[below,font=\scriptsize] {clash: $\Top>f$} (y1);
\end{tikzpicture}
\caption{The nonterminal cross event of \cref{ex:cross-event}, for
$\phi=\{x\le f(a,y),\ y\le f(a,b)\}$ with $c=(1,0)$, the query $x\le y$,
and $w=11$. The comparison $(0,1)$ is cross-admitted through
$y^+\in U_1\cap L_0$. The bridge $D(P_1,Y_0)$ descends one matching letter to
$D(P_2,Y_1)$, where top meets an upper constructor bound.}
\Description{The word consisting of two ones has cuts zero, one and two,
with its one-letter suffix at cut one matching the prefix at cut zero.
A middle row connects P one through the shared variable y positive to Y zero.
Descending child one on both sides gives a closure pair from P two, labeled
top, to Y one, bounded above by a constructor. These labels clash.}
\label{fig:cross-event}
\end{figure}


\begin{example}[A self event]\label{ex:self-event}
Keep the arrow variance $c=(1,0)$ and take $\phi=\{x\le f(x,b)\}$, that is,
$x\le x\to b$, with the query $x\le y$ and $y$ unconstrained. The signed
literals are $x^+\le f(x^-,b^+)$ and $f(x^+,b^-)\le x^-$, and the closure is
the identity. Along $w=0$, the upper sets are $U_0=\{x^+\}$ and
$U_1=\{x^-\}$, and the lower sets are $L_0=\{y^+\}$ and $L_1=\varnothing$.
There are no constant literals and $U_j\cap L_j=\varnothing$, so neither cut
is ready, and $\childtest(U_1)$ fails because $x^-$ roots no upper constructor
bound. The terminal comparison $(0,1)$ is not cross-admitted, since
$L_0\cap U_1=\varnothing$, but it is self-admitted:
$U_0\cap\flip U_1=\{x^+\}$. The self event alone shows that $w=0$ is not
left unsafe. Semantically, $x\le_c x\to b$ excludes $x=\top$, so a top prefix
of $x^+$ along $0$ needs a bottom domain leaf, $x=f(\bot,t)$; but then
$x\le_c x\to b$ forces $x\le_c\bot$, which is impossible.
In the extension, $P_0=x^+$, $Q_0=x^-$, $P_1=\top$, and $Q_1=\bot$, with the
new links $f(P_1,B^0)\le x^+$ and $x^-\le f(Q_1,T^1)$. Decomposing the old
lower literal $f(x^+,b^-)\le x^-$ against the new upper link gives
$D(P_0,Q_1)$, and the new lower constructor bound on $P_0$ clashes with
$Q_1=\bot$. This is the terminal self case of \cref{thm:clashes}.
\end{example}
